\section{Wholes And Parts}

Counting begins when a distinction can be carried. A difference that
appears once and vanishes leaves no count behind. A difference that can
be carried into a next reading gives the grammar something to step with.
The step is not yet arithmetic in its finished form. It is the more
primitive ability to say that a reading has continued, or that a reading
has split. The count is born at the fork between those two possibilities.

The chapter calls the two possibilities same-parity and
different-parity readings. In a same-parity reading, the grammar carries
the current distinction as a whole. The new reading belongs with the old
one. It extends the record without breaking the unit that the record has
been treating as one. This is counting by wholes. A tally mark added to
a row, a repeated beat in a metronome, a trial that agrees with the
previous trial, or a sample that remains inside the same tolerance all
illustrate the same permission: the grammar reads continuation.

In a different-parity reading, the grammar does not carry the prior
whole intact. It reads a change of side. The old unit is no longer
enough; the record must expose a part, a contrast, a residue, or a
branch. This is counting by parts. A crack in a ceramic plate, a bit
flip in a memory cell, a phase reversal in a signal, or a failed match
in a repeated trial all illustrate the same permission: the grammar
reads fracture.

Neither reading is secondary. If the grammar could only count wholes, it
would have no way to notice when a whole had ceased to hold. If it could
only count parts, it would have no way to preserve a unit long enough to
measure it. Measurement needs both. It must be able to pool readings
that belong together, and it must be able to select the boundary at which
belonging fails. The whole-reading and the part-reading are not rival
theories. They are conjugate obligations of the first distinction.

This is why the grammar treats count as already oriented. A count is not
merely an accumulation of strokes. It is a decision about what the next
stroke is allowed to carry. When a laboratory repeats a calibration, an
agreement is not just another mark; it is a mark that preserves the
calibration as a whole. When a reading drifts outside tolerance, the new
mark does not merely increase the total; it forces a part to appear, a
difference inside what had been treated as a stable unit. The ledger must
know which kind of mark it has accepted.

The distinction between wholes and parts also explains why number cannot
be separated from order at the foundation. To count a whole is to say
that the next reading is ordered with the previous one by preservation.
To count a part is to say that the next reading is ordered by contrast.
The grammar does not need a completed line of natural numbers in order
to begin. It needs a permitted successor act, and the successor act must
say whether it carries the unit forward or exposes a difference within
it.

This also prevents a common mistake about discreteness. The primitive
count is not a heap of isolated pellets. It is a disciplined record of
what can be repeated and what must be split. A sequence of identical
marks is meaningful only because the grammar has permitted them to be
read as the same kind of mark. A sequence of alternating marks is
meaningful only because the grammar has permitted the alternation to be
read as a difference that remains recoverable. In both cases the count
depends on a carried reading, not on a preexisting pile of units.

The result is the first mathematical shape of the book. The grammar
permits a number to begin as a record of permitted continuation and
permitted fracture. It forbids the number from pretending to be a
completed infinite object. It identifies no universal identity yet. It
only carries enough same-reading and different-reading structure to ask
for a next refinement. That next refinement will be finite, and its
unfinished remainder will matter.
